% GENERATED FILE. Edit canonical lessons, then run npm run generate:books.
% canonical-source-hash: fnv1a64:516d75a8042c0d2c
% canonical-lessons: GU-C478-kaph, GU-C478-udhras, GU-C478-kharas, GU-C478-khansvu, GU-C478-kabjiyat, GU-R478-review-471-478

\chapter{A Cough That Lingers}
\label{ch:gu-a-cough-that-lingers}
\hlchaptermodality{478}

\begin{chapteropening}
\textbf{By the end of this chapter:} \emph{I can describe a cough and a sore throat.}

Complete the last lesson of chapter 478: I can describe a cough and a sore throat.
\end{chapteropening}

\section[\texorpdfstring{\gu{કફ} (kaph)}{કફ (kaph)}]{\gu{કફ} (kaph) — phlegm}

\label{lesson:GU-C478-kaph}



\begin{quote}
Before the new one: say the Gujarati for to throw up, then the Gujarati for to stomach.
\end{quote}

\subsection*{You'll want to know: \gu{કફ}}
\textbf{\gu{કફ}} — \emph{kaph} — \textquotedblleft{}phlegm\textquotedblright{}.

\begin{quote}
\gu{શરદીને} \gu{લીધે} \gu{મારી} \gu{છાતીમાં} \gu{કફ} \gu{છે}.
\end{quote}

\begin{quote}
\emph{śardīne līdhe mārī chātīmā̃ kaph che.}
\end{quote}

\begin{quote}
My chest is full of phlegm because of the cold.
\end{quote}

A masculine noun.

\begin{cousinweb}
\textbf{\gu{કફ}} (\emph{kaph}) is a tatsama, Sanskrit \emph{kapha}, \textquotedblleft{}phlegm\textquotedblright{}, one of the three humours of the old medicine, as Monier-Williams records it.
\end{cousinweb}

\begin{grammarlens}[title={What you've built}]
One more word for describing what you lived and felt.
\end{grammarlens}

\subsection*{Your turn}
\begin{itemize}
\raggedright
  \item \emph{Say it:} \emph{kaph}
  \item \emph{Say it:} \emph{kaph}, once more
  \item \emph{Say it:} \emph{śardīne līdhe mārī chātīmā̃ kaph che.} — My chest is full of phlegm because of the cold.
  \item \emph{Recall:} say the Gujarati for to throw up, then the Gujarati for to stomach, then say \emph{kaph} again
\end{itemize}

\begin{tcolorbox}[breakable,colback=teal!4,colframe=teal!35!black,title={Before you move on}]
What does \gu{કફ} mean? (\textquotedblleft{}Phlegm\textquotedblright{}.) Say it once more.
\end{tcolorbox}
\section[\texorpdfstring{\gu{ઉધરસ} (udhras)}{ઉધરસ (udhras)}]{\gu{ઉધરસ} (udhras) — a lingering cough; a cough}

\label{lesson:GU-C478-udhras}



\begin{quote}
Before the new one: say the Gujarati for to stomach, then the Gujarati for phlegm.
\end{quote}

\subsection*{You'll want to know: \gu{ઉધરસ}}
\textbf{\gu{ઉધરસ}} — \emph{udhras} — \textquotedblleft{}a lingering cough; a cough\textquotedblright{}.

\begin{quote}
\gu{દાદીને} \gu{આખી} \gu{રાત} \gu{ઉધરસ} \gu{આવી}.
\end{quote}

\begin{quote}
\emph{dādīne ākhī rāt udhras āvī.}
\end{quote}

\begin{quote}
Grandmother coughed all night.
\end{quote}

A feminine noun.

\begin{cousinweb}
\textbf{\gu{ઉધરસ}} (\emph{udhras}) goes back to a reconstructed \emph{*uddhrasā-}, \textquotedblleft{}coughing up\textquotedblright{}; CDIAL lists the Gujarati word, \textquotedblleft{}cough\textquotedblright{} (CDIAL 2030).
\end{cousinweb}

\begin{grammarlens}[title={What you've built}]
One more word for describing what you lived and felt.
\end{grammarlens}

\subsection*{Your turn}
\begin{itemize}
\raggedright
  \item \emph{Say it:} \emph{udhras}
  \item \emph{Say it:} \emph{udhras}, once more
  \item \emph{Say it:} \emph{dādīne ākhī rāt udhras āvī.} — Grandmother coughed all night.
  \item \emph{Recall:} say the Gujarati for to stomach, then the Gujarati for phlegm, then say \emph{udhras} again
\end{itemize}

\begin{tcolorbox}[breakable,colback=teal!4,colframe=teal!35!black,title={Before you move on}]
What does \gu{ઉધરસ} mean? (\textquotedblleft{}A lingering cough; a cough\textquotedblright{}.) Say it once more.
\end{tcolorbox}
\section[\texorpdfstring{\gu{ખરાશ} (kharāś)}{ખરાશ (kharāś)}]{\gu{ખરાશ} (kharāś) — a scratchy throat}

\label{lesson:GU-C478-kharas}



\begin{quote}
Before the new one: say the Gujarati for phlegm, then the Gujarati for a lingering cough.
\end{quote}

\subsection*{You'll want to know: \gu{ખરાશ}}
\textbf{\gu{ખરાશ}} — \emph{kharāś} — \textquotedblleft{}a scratchy throat\textquotedblright{}.

\begin{quote}
\gu{ઠંડું} \gu{પાણી} \gu{પીધા} \gu{પછી} \gu{ગળામાં} \gu{ખરાશ} \gu{થઈ}.
\end{quote}

\begin{quote}
\emph{ṭhanḍũ pāṇī pīdhā pachī gaḷāmā̃ kharāś thaī.}
\end{quote}

\begin{quote}
After drinking cold water I got a scratchy throat.
\end{quote}

A feminine noun.

\begin{cousinweb}
\textbf{\gu{ખરાશ}} (\emph{kharāś}) is usually traced to Persian \emph{xarāš}, \textquotedblleft{}a scratch, a scraping\textquotedblright{}.
\end{cousinweb}

\begin{grammarlens}[title={What you've built}]
One more word for describing what you lived and felt.
\end{grammarlens}

\subsection*{Your turn}
\begin{itemize}
\raggedright
  \item \emph{Say it:} \emph{kharāś}
  \item \emph{Say it:} \emph{kharāś}, once more
  \item \emph{Say it:} \emph{ṭhanḍũ pāṇī pīdhā pachī gaḷāmā̃ kharāś thaī.} — After drinking cold water I got a scratchy throat.
  \item \emph{Recall:} say the Gujarati for phlegm, then the Gujarati for a lingering cough, then say \emph{kharāś} again
\end{itemize}

\begin{tcolorbox}[breakable,colback=teal!4,colframe=teal!35!black,title={Before you move on}]
What does \gu{ખરાશ} mean? (\textquotedblleft{}A scratchy throat\textquotedblright{}.) Say it once more.
\end{tcolorbox}
\section[\texorpdfstring{\gu{ખાંસવું} (khā̃svũ)}{ખાંસવું (khā̃svũ)}]{\gu{ખાંસવું} (khā̃svũ) — to cough}

\label{lesson:GU-C478-khansvu}



\begin{quote}
Before the new one: say the Gujarati for a lingering cough, then the Gujarati for a scratchy throat.
\end{quote}

\subsection*{You'll want to know: \gu{ખાંસવું}}
\textbf{\gu{ખાંસવું}} — \emph{khā̃svũ} — \textquotedblleft{}to cough\textquotedblright{}.

\begin{quote}
\gu{દર્દી} \gu{આખી} \gu{રાત} \gu{ખાંસતો} \gu{હતો}.
\end{quote}

\begin{quote}
\emph{dardī ākhī rāt khā̃sto hato.}
\end{quote}

\begin{quote}
The patient was coughing all night.
\end{quote}

Past (he …): \gu{તે} \gu{આખી} \gu{રાત} \gu{ખાંસ્યો}, \emph{te ākhī rāt khā̃syo}.

\begin{cousinweb}
\textbf{\gu{ખાંસવું}} (\emph{khā̃svũ}) is inherited from Sanskrit \emph{kāsatē}, \textquotedblleft{}coughs\textquotedblright{}; CDIAL lists the Gujarati verb (CDIAL 3135), beside the noun \textbf{\gu{ખાંસી}} (\emph{khā̃sī}), \textquotedblleft{}a cough\textquotedblright{}.
\end{cousinweb}

\begin{grammarlens}[title={What you've built}]
One more word for describing what you lived and felt.
\end{grammarlens}

\subsection*{Your turn}
\begin{itemize}
\raggedright
  \item \emph{Say it:} \emph{khā̃svũ}
  \item \emph{Say it:} \emph{khā̃svũ}, once more
  \item \emph{Say it:} \emph{dardī ākhī rāt khā̃sto hato.} — The patient was coughing all night.
  \item \emph{Recall:} say the Gujarati for a lingering cough, then the Gujarati for a scratchy throat, then say \emph{khā̃svũ} again
\end{itemize}

\begin{tcolorbox}[breakable,colback=teal!4,colframe=teal!35!black,title={Before you move on}]
What does \gu{ખાંસવું} mean? (\textquotedblleft{}To cough\textquotedblright{}.) Say it once more.
\end{tcolorbox}
\section[\texorpdfstring{\gu{કબજિયાત} (kabjiyāt)}{કબજિયાત (kabjiyāt)}]{\gu{કબજિયાત} (kabjiyāt) — constipation}

\label{lesson:GU-C478-kabjiyat}



\begin{quote}
Before the new one: say the Gujarati for a scratchy throat, then the Gujarati for to cough.
\end{quote}

\subsection*{You'll want to know: \gu{કબજિયાત}}
\textbf{\gu{કબજિયાત}} — \emph{kabjiyāt} — \textquotedblleft{}constipation\textquotedblright{}.

\begin{quote}
\gu{દાક્તરે} \gu{કહ્યું} \gu{કે} \gu{કબજિયાત} \gu{માટે} \gu{વધારે} \gu{પાણી} \gu{પીઓ}.
\end{quote}

\begin{quote}
\emph{dāktare kahyũ ke kabjiyāt māṭe vadhāre pāṇī pīo.}
\end{quote}

\begin{quote}
The doctor said to drink more water for constipation.
\end{quote}

A feminine noun.

\begin{cousinweb}
\textbf{\gu{કબજિયાત}} (\emph{kabjiyāt}) is usually traced to Arabic \emph{qabḍ}, \textquotedblleft{}contraction, seizing\textquotedblright{}, through Persian, with an abstract ending.
\end{cousinweb}

\begin{grammarlens}[title={What you've built}]
One more word for describing what you lived and felt.
\end{grammarlens}

\subsection*{Your turn}
\begin{itemize}
\raggedright
  \item \emph{Say it:} \emph{kabjiyāt}
  \item \emph{Say it:} \emph{kabjiyāt}, once more
  \item \emph{Say it:} \emph{dāktare kahyũ ke kabjiyāt māṭe vadhāre pāṇī pīo.} — The doctor said to drink more water for constipation.
  \item \emph{Recall:} say the Gujarati for a scratchy throat, then the Gujarati for to cough, then say \emph{kabjiyāt} again
\end{itemize}

\begin{tcolorbox}[breakable,colback=teal!4,colframe=teal!35!black,title={Before you move on}]
What does \gu{કબજિયાત} mean? (\textquotedblleft{}Constipation\textquotedblright{}.) Say it once more.
\end{tcolorbox}
\section[(dialogue)]{Looking back: from How People Look to A Cough That Lingers}

\label{lesson:GU-R478-review-471-478}



\begin{quote}
Say the Gujarati for complexion and height.
\end{quote}

\subsection*{Your turn}
Say these aloud:

\begin{itemize}
\raggedright
  \item complexion, height, the body, wrinkles, grey hairs
  \item a vein, a muscle, the skull, an intestine, the stomach
  \item a thigh, the sole of the foot, the navel, the calf of the leg, an armpit
  \item saliva, to chew, to suck, to spit, the palate
  \item blind, deaf, silent, eyesight, to touch
  \item to snore, slumber, staying up all night, to lie down to sleep, sleeplessness
  \item a belch, bloating, nausea, to throw up, to stomach
  \item phlegm, a lingering cough, a scratchy throat, to cough, constipation
\end{itemize}

\begin{tcolorbox}[breakable,colback=teal!4,colframe=teal!35!black,title={Before you move on}]
Complexion? (\textbf{\gu{વાન}}, \emph{vān}.) Height? (\textbf{\gu{કદ}}, \emph{kad}.) The body? (\textbf{\gu{દેહ}}, \emph{deh}.) Wrinkles? (\textbf{\gu{વળિયાં}}, \emph{vaḷiyā̃}.) Grey hairs? (\textbf{\gu{પળિયાં}}, \emph{paḷiyā̃}.) A vein? (\textbf{\gu{નસ}}, \emph{nas}.) A muscle? (\textbf{\gu{સ્નાયુ}}, \emph{snāyu}.) The skull? (\textbf{\gu{ખોપરી}}, \emph{khoprī}.) An intestine? (\textbf{\gu{આંતરડું}}, \emph{ā̃tarḍũ}.) The stomach? (\textbf{\gu{હોજરી}}, \emph{hojrī}.) A thigh? (\textbf{\gu{જાંઘ}}, \emph{jā̃gh}.) The sole of the foot? (\textbf{\gu{તળિયું}}, \emph{taḷiyũ}.) The navel? (\textbf{\gu{ડૂંટી}}, \emph{ḍū̃ṭī}.) The calf of the leg? (\textbf{\gu{પિંડી}}, \emph{pinḍī}.) An armpit? (\textbf{\gu{બગલ}}, \emph{bagal}.) Saliva? (\textbf{\gu{લાળ}}, \emph{lāḷ}.) To chew? (\textbf{\gu{ચાવવું}}, \emph{cāvvũ}.) To suck? (\textbf{\gu{ચૂસવું}}, \emph{cūsvũ}.) To spit? (\textbf{\gu{થૂંકવું}}, \emph{thū̃kvũ}.) The palate? (\textbf{\gu{તાળવું}}, \emph{tāḷvũ}.) Blind? (\textbf{\gu{અંધ}}, \emph{andh}.) Deaf? (\textbf{\gu{બહેરું}}, \emph{baherũ}.) Silent? (\textbf{\gu{મૂંગું}}, \emph{mū̃gũ}.) Eyesight? (\textbf{\gu{નજર}}, \emph{najar}.) To touch? (\textbf{\gu{અડકવું}}, \emph{aḍakvũ}.) To snore? (\textbf{\gu{ઘોરવું}}, \emph{ghorvũ}.) Slumber? (\textbf{\gu{નીંદર}}, \emph{nī̃dar}.) Staying up all night? (\textbf{\gu{જાગરણ}}, \emph{jāgraṇ}.) To lie down to sleep? (\textbf{\gu{પોઢવું}}, \emph{poḍhvũ}.) Sleeplessness? (\textbf{\gu{અનિદ્રા}}, \emph{anidrā}.) A belch? (\textbf{\gu{ઓડકાર}}, \emph{oḍkār}.) Bloating? (\textbf{\gu{આફરો}}, \emph{āphro}.) Nausea? (\textbf{\gu{ઉબકો}}, \emph{ubko}.) To throw up? (\textbf{\gu{ઓકવું}}, \emph{okvũ}.) To stomach? (\textbf{\gu{જીરવવું}}, \emph{jīravvũ}.) Phlegm? (\textbf{\gu{કફ}}, \emph{kaph}.) A lingering cough? (\textbf{\gu{ઉધરસ}}, \emph{udhras}.) A scratchy throat? (\textbf{\gu{ખરાશ}}, \emph{kharāś}.) To cough? (\textbf{\gu{ખાંસવું}}, \emph{khā̃svũ}.) Constipation? (\textbf{\gu{કબજિયાત}}, \emph{kabjiyāt}.)
\end{tcolorbox}
